% sec_front.tex -- chunk A: title, abstract, \maketitle, Section 1 (Introduction).
% Source comments (% src: ...) point to the repository and are removed before submission.
% Labels defined here: sec:intro, sec:setting, sec:results, sec:context, sec:strategy, sec:organisation,
%   eq:sigma, eq:jacobi, eq:cp, eq:law, eq:paper1-averaged, eq:osc, eq:smoothing, eq:rho, eq:strip, eq:epsLambda,
%   eq:thetap, def:sv, def:smoothing, def:A, def:rate, thm:main, thm:pure, thm:shift, thm:engine, rem:main-*, rem:pure-*,
%   rem:shift-*, rem:engine-scope, fig:dependency.
% Labels used here but defined in other chunks (see NOTES.md, "Cross-chunk labels"): sec:prelim, sec:tools,
%   sec:shift, sec:smoothing, sec:frame, sec:outer, sec:hardedge, sec:extraction, lem:potter, lem:smoothing,
%   lem:saturation, prop:stability, prop:equilibrium, prop:frame, rem:quotient-p1, prop:UR-comparison,
%   rem:alpha-ge-1, rem:base-window, app:L2, app:resolvent, app:apriori; sec:numerics, sec:further (chunk A, sec_back).

\providecommand{\Ath}[1]{\textup{(A$_{#1}$)}}
\providecommand{\ceil}[1]{\lceil #1\rceil}

\title[Recurrence coefficients with slowly varying weights]{Recurrence coefficients of discrete measures with
power-law atoms and slowly varying weights}
\author{Gordon [surname] (draft)}
\date{Draft of 14 October 2026}

\begin{abstract}
Let $\sigma=\sum_{k\ge1}k^{-s}\ell(k)\,\delta_{k^{-p}}$ with $p>0$, $s>1$ and $\ell$ slowly varying; its Jacobi
matrix is a compact operator with eigenvalues $k^{-p}$. We prove that the recurrence coefficients satisfy
$n^pb_n\to c(p)$ and $n^pa_n\to2c(p)$, where
$c(p)=\frac14\bigl(\sqrt\pi\,\Gamma(1/p+1)/\Gamma(1/p+\frac12)\bigr)^p$ is the constant of the constrained
equilibrium problem solved in a companion paper. This holds for every slowly varying $\ell$, including unbounded ones such as
$\exp\sqrt{\log k}$. If $\ell$ extends analytically to a sector, the error is
$O(1/n)$; for pure powers this gives $n^pb_n=c(p)(1+O(1/n))$ and $a_n/(2b_n)=1+O(1/n)$. The proof has three parts. A
Riemann--Hilbert analysis of discrete orthogonal polynomials with nodes accumulating at a hard edge controls the
non-analytic factor by a Conway--Deift comparison with a Bessel model instead of a local parametrix. A smoothing step
reduces arbitrary to analytic slowly varying factors. The Christoffel transformation $\sigma\mapsto x\sigma$ shifts
$s$ by $p$ and preserves the law together with its rate. In a further paper we show that bounded but rough weights
can destroy the pointwise law.
\end{abstract}
% Word count of the abstract (pdftotext of the compiled page): about 200.

\maketitle

\section{Introduction}\label{sec:intro}

\subsection{Setting}\label{sec:setting}
% src: OUTLINE.md §1.3 (story); papers/paper1/main.tex §1.
A compact Jacobi operator is determined by its spectral measure. If its eigenvalues are $k^{-p}$ and its spectral
weights decay like $k^{-s}$, how do its matrix entries behave? This paper answers the question pointwise, for every
$p>0$ and $s>1$, and shows that the answer is stable when the weights are multiplied by an arbitrary slowly varying
factor.

Concretely, for $p>0$, $s>1$ and a sequence $\ell\colon\N\to(0,\infty)$ we study the discrete measure
\begin{equation}\label{eq:sigma}
  \sigma_{p,s,\ell}:=\sum_{k\ge1}k^{-s}\ell(k)\,\delta_{k^{-p}},\qquad \sigma_{p,s}:=\sigma_{p,s,1}.
\end{equation}
It lives on $(0,1]$ and has infinitely many atoms, which accumulate only at $0$. Let $P_n$ be its monic orthogonal
polynomials and $h_n:=\int P_n^2\diff\sigma$. They satisfy the three-term recurrence
\begin{equation}\label{eq:jacobi}
  xP_n=P_{n+1}+a_nP_n+b_n^2P_{n-1}\qquad(n\ge0,\ b_0:=0,\ b_n>0),
\end{equation}
so that $b_n^2=h_n/h_{n-1}$. The numbers $a_n$ and $b_n$ are the diagonal and off-diagonal entries of the Jacobi
matrix of $\sigma$. (The companion paper \cite{Paper1} writes $a_{n+1}$ for our $a_n$; the $b_n$ agree. See
Section~\ref{sec:prelim}.)

The constants that govern the answer come from \cite{Paper1}:
\begin{equation}\label{eq:cp}
  I(p):=\frac{\sqrt\pi}{2}\,\frac{\Gamma(\frac1p+\frac12)}{\Gamma(\frac1p+1)},\qquad
  c(p):=\frac14\Bigl(\frac{\pi}{2I(p)}\Bigr)^{p}
  =\frac14\Bigl(\frac{\sqrt\pi\,\Gamma(\frac1p+1)}{\Gamma(\frac1p+\frac12)}\Bigr)^{p}.
\end{equation}
For example $c(1)=\frac12$ and $c(2)=\pi^2/16$. The question is whether
\begin{equation}\label{eq:law}
  n^pb_n\longrightarrow c(p)\qquad\text{and}\qquad n^pa_n\longrightarrow2c(p)\qquad(n\to\infty).
\end{equation}
We call \eqref{eq:law} the \emph{pointwise law}.

\emph{The averaged law.} The companion paper \cite{Paper1} proves \eqref{eq:law} on average. Its
Corollary~1.2, specialised to the atoms $x_k=k^{-p}$, states: if the weights $w_k\ge0$ satisfy $\sum_kw_k<\infty$ and
$w_k\ge c_1k^{-s'}$ for some constants $c_1>0$ and $s'\in\R$, then the recurrence coefficients of
$\sum_kw_k\delta_{k^{-p}}$ satisfy
\begin{equation}\label{eq:paper1-averaged}
  \frac{2}{n^2}\sum_{j=1}^{n-1}(n-j)\log\bigl(j^pb_j\bigr)\longrightarrow\log c(p).
\end{equation}
The constant arises from an explicitly solvable constrained equilibrium problem \cite[Theorem~1.3]{Paper1}. Its
solution $\mu^*$ is an average of arcsine laws, and it is exactly the zero distribution that a Jacobi matrix with
entries $a_n\approx2b_n\approx2c\,n^{-p}$ would produce; the saturation condition of the equilibrium problem fixes
$c=c(p)$ \cite[\S1.2]{Paper1}. So $c(p)$ is the value predicted for the pointwise law.

\emph{Why the pointwise law is subtle.} The statement \eqref{eq:paper1-averaged} is a double Ces\`aro mean of
$\log(j^pb_j)$. Passing from it to \eqref{eq:law} is a Tauberian step, and it needs a regularity of the sequence
$j^pb_j$ that the averaged argument does not provide. That regularity must come from the weights. Indeed, the averaged
law holds under the one-sided bound $w_k\ge c_1k^{-s'}$ alone, while the pointwise law can fail for weights that
oscillate within a fixed factor of $k^{-s}$ (see the announcement in \S\ref{sec:context}). This paper shows that slow
variation of the weights is enough.

\subsection{Results}\label{sec:results}
We use the following form of slow variation, which needs no measurability or representation theorem.

\begin{definition}[slow variation]\label{def:sv}
% src: research/G1/PROOF.md §3.14 (sequence form; §3.8/§3.9(i)).
A sequence $\ell\colon\N\to(0,\infty)$ is \emph{slowly varying} if, for every $\lambda\ge1$,
\begin{equation}\label{eq:osc}
  \osc_n(\lambda):=\max\bigl\{\,|\log\ell(k)-\log\ell(n)|:\ k\in\N,\ n/\lambda\le k\le\lambda n\,\bigr\}
  \longrightarrow0\qquad(n\to\infty).
\end{equation}
\end{definition}

Equivalently, $\ell(k)/\ell(n)\to1$ uniformly for $k\in[n/\lambda,\lambda n]$, for every fixed $\lambda$; the case
$\lambda=2$ already implies the others (Section~\ref{sec:prelim}). Examples are $(\log(k+1))^\beta$ with
$\beta\in\R$, $\exp((\log k)^\theta)$ with $0<\theta<1$, and $\exp(\sqrt{\log k}\,\sin\log\log k)$ (for $k\ge3$, say).
Neither boundedness nor monotonicity is required.

The rate in our main theorem is measured by how well $\ell$ can be smoothed.

\begin{definition}[smoothing error]\label{def:smoothing}
% src: research/G1/RH/route2/ROUTE2.md §1 (RedA).
Let $L_\ell\colon\R\to\R$ be the piecewise-linear interpolant of the points $(\log k,\log\ell(k))$, $k\ge1$, extended
by the constant $\log\ell(1)$ on $(-\infty,0]$. Let $G(y):=e^{-y^2/2}/\sqrt{2\pi}$ and
\begin{equation}\label{eq:smoothing}
  \Lambda_\ell:=L_\ell*G,\qquad
  \rho_\ell(m):=\sup_{k\ge m}\bigl|\log\ell(k)-\Lambda_\ell(\log k)\bigr|\qquad(m\ge1).
\end{equation}
\end{definition}

The function $\Lambda_\ell$ is entire, and $\rho_\ell(m)\to0$ as $m\to\infty$ whenever $\ell$ is slowly varying
(Lemma~\ref{lem:smoothing}).

For the sharp rate we need an analytic extension of $\ell$ to a sector, written in the logarithmic variable
$\xi=\log t$. For $\vartheta>0$ and $c_0>0$ put $S_\vartheta:=\{\xi\in\C:\ \Re\xi\ge-c_0,\ |\Im\xi|\le\vartheta\}$.

\begin{definition}[analytic slowly varying factors]\label{def:A}
% src: (H_an) ROUTE2.md §2 (p > 1, Rule R F4 domain); (H_an^{x1}) ROUTE2_xframe1.md §2 (p = 1);
% (H_an^{<1}) ROUTE2_plt1.md §1 (p < 1). Not unified (OUTLINE §7 R9).
Let $\vartheta>0$. A sequence $\ell$ satisfies \Ath{\vartheta} if $\ell(k)=\exp\Lambda(\log k)$ for all $k\ge1$,
where, for some $c_0>0$, the function $\Lambda$ is analytic on $S_\vartheta$, real on $S_\vartheta\cap\R$, bounded on
$S_\vartheta\cap\{\Re\xi\le0\}$, and
\begin{equation}\label{eq:epsLambda}
  \varepsilon_\Lambda(T):=\sup\bigl\{|\Lambda'(\xi)|:\ \xi\in S_\vartheta,\ \Re\xi\ge T\bigr\}\longrightarrow0
  \qquad(T\to\infty).
\end{equation}
We use the half-width
\begin{equation}\label{eq:thetap}
  \vartheta_p:=\pi/p\quad\text{for }p>1,\qquad\vartheta_p:=\pi/2\quad\text{for }0<p\le1.
\end{equation}
\end{definition}

In words, \Ath{\vartheta_p} asks that $\ell$ extend to an analytic function $\ell(t)=e^{\Lambda(\log t)}$ on the sector
$|\arg t|\le\vartheta_p$ (for large $|t|$, and bounded near $t=1$), with $t\ell'(t)/\ell(t)\to0$ there. A sequence
satisfying \Ath{\vartheta} is slowly varying, since $|\Lambda(\log k)-\Lambda(\log n)|\le\varepsilon_\Lambda(\log
n-\log\lambda)\log\lambda$ for $n/\lambda\le k\le\lambda n$. The sequence $\ell(k)=\exp\sqrt{\log(k+e)}$ satisfies
\Ath{\vartheta_p} for every $p>0$, because the branch points of $\Lambda(\xi)=\sqrt{\log(e^\xi+e)}$ lie on the lines
$|\Im\xi|=\pi$. The smoothed function $\Lambda_\ell$ of any slowly varying $\ell$ satisfies \Ath{\vartheta} for every
$\vartheta>0$ (Lemma~\ref{lem:smoothing}). For $p>1$ a further example is
$\Lambda(\xi)=(\log(e^\xi+e))^\theta$ with $\theta<1$.
% (The theta-example is stated for p > 1 only: the repository checks it there (ROUTE2 §2); OUTLINE §2.)

Finally, $M_0=M_0(p)\ge4$ denotes an integer that depends only on $p$; it is fixed in
Section~\ref{sec:smoothing} (Lemma~\ref{lem:saturation}) and marks the indices $k<n/M_0$ at which the spectral weights
of $P_n$ are exponentially saturated.
% src: PROOF.md §3.14 notation line, §3.3: M_0 is the least integer >= 4 with g_p^+ >= 1 on (0, 2/M_0].

\begin{maintheorem}[main theorem]\label{thm:main}
% src: THEOREMS.md §2.8; PROOF.md §3.21 Cor. B_all (p > 1), §3.20 Cor. B_{p<1}, §3.23 Cor. B_{p=1}^{all}
% (with §3.19 B_{p=1}^+); rate (ii) in the re-ordered form (OUTLINE §0 item 5).
Let $p>0$, $s>1$, and let $\ell\colon\N\to(0,\infty)$ be slowly varying. Let
\[
  \sigma=\sigma_{p,s,\ell}:=\sum_{k\ge1}k^{-s}\ell(k)\,\delta_{k^{-p}}.
\]
Then
\begin{enumerate}
\item[\textup{(i)}] $n^pb_n\to c(p)$ and $n^pa_n\to2c(p)$;
\item[\textup{(ii)}] more precisely, $|n^pb_n-c(p)|+|n^pa_n-2c(p)|=O\bigl(1/n+\rho_\ell(\ceil{n/M_0})\bigr)$;
\item[\textup{(iii)}] if $\ell$ satisfies \Ath{\vartheta_p}, then $n^pb_n=c(p)(1+O(1/n))$ and
  $a_n/(2b_n)=1+O(1/n)$.
\end{enumerate}
Here $c(p)=\frac14(\pi/(2I(p)))^p=\frac14(\sqrt\pi\,\Gamma(1/p+1)/\Gamma(1/p+\frac12))^p$, with
$I(p)=(\sqrt\pi/2)\Gamma(1/p+\frac12)/\Gamma(1/p+1)$. For example, $c(1)=\frac12$ and $c(2)=\pi^2/16$.
\end{maintheorem}

\begin{remark}\label{rem:main-finite}
The measure is finite: by the Potter bounds (Lemma~\ref{lem:potter}), $\ell(k)\le C_\eta k^\eta$ for every $\eta>0$,
so $\sum_kk^{-s}\ell(k)<\infty$. Thus $\sigma$ is a finite measure on $(0,1]$ with infinitely many atoms, and its
Jacobi matrix is a compact operator with eigenvalues $k^{-p}$.
\end{remark}

\begin{remark}\label{rem:main-unbounded}
No boundedness of $\ell$ is assumed. For example, the weights
\[
  k^{-s}\exp\sqrt{\log k}\qquad\text{and}\qquad k^{-s}\exp\bigl(\sqrt{\log k}\,\sin\log\log k\bigr)\quad(k\ge3)
\]
are covered. The second one is not even eventually comparable with any fixed power of $\log k$.
\end{remark}

\begin{remark}\label{rem:main-rate}
Part (iii) bounds the deviation from $c(p)$ at order $1/n$, but it does not identify the $1/n$ term. In particular we
do not claim that the slowly varying factor changes $b_n$ only by a relative $o(1/n)$, and we do not determine its
effect at order $1/n$. A conjecture, supported by numerics, is formulated in \S\ref{sec:further}.
\end{remark}

\begin{remark}\label{rem:main-averaged}
The averaged law \eqref{eq:paper1-averaged} holds for much rougher weights. Theorem~\ref{thm:main} is the pointwise
statement, and \S\ref{sec:context} explains why some regularity of the weights is necessary for it.
\end{remark}

For pure powers the rate is the sharp one.

\begin{maintheorem}[pure powers, with rate]\label{thm:pure}
% src: THEOREMS.md H2 (RH_1(r), RH_p, RH_p seam). In the paper: Theorem 4 at Lambda = 0, plus Theorem 3.
For every $p>0$ and $s>1$ the measure $\sigma_{p,s}:=\sum_{k\ge1}k^{-s}\delta_{k^{-p}}$ satisfies
\[
  n^pb_n=c(p)\bigl(1+O(1/n)\bigr),\qquad a_n/(2b_n)=1+O(1/n).
\]
Equivalently, $|n^pb_n-c(p)|+|n^pa_n-2c(p)|=O(1/n)$.
\end{maintheorem}

\begin{remark}[a classical instance]\label{rem:pure-lommel}
% src: novelty_G1RH.md §1.3 (Van Assche math/9510214 §3, quoted there); checks/lommel_check.py, .out.
For $(p,s)=(2,2)$ the coefficients are explicit, and Theorem~\ref{thm:pure} can be checked directly. Van Assche
\cite[\S3]{VA} records that the Lommel polynomials $h_{n,\nu}$ satisfy $h_{n+1,\nu}(x)=2x(n+\nu)h_{n,\nu}(x)-h_{n-1,\nu}(x)$,
that their orthonormal versions have the coefficients $a_n=0$ and $\lambda_n=[4(n+\nu)(n+\nu-1)]^{-1}$, and that the
spectrum is the closure of $\{1/j_{k,\nu-1}\}$, the $j_{k,\nu-1}$ being the zeros of $J_{\nu-1}$. For $\nu=\frac32$
these zeros are $\pm\pi k$. Passing from this symmetric measure to its image under $y\mapsto\pi^2y^2$ (the usual
quadratic transformation) leads to
\[
  b_n=\frac{\pi^2}{(4n+1)\sqrt{(4n-1)(4n+3)}},\qquad a_n=\frac{2\pi^2}{(4n+1)(4n+5)}\qquad(n\ge1)
\]
for $\sigma_{2,2}$. We have confirmed numerically that these values agree with the coefficients of $\sigma_{2,2}$
computed from the Hankel determinants of its moments $\zeta(2j+2)$ (to $90$ digits for $n\le9$); we have not checked
the identification of the masses against a source. They give $n^2b_n/c(2)=1-\frac1{2n}+O(n^{-2})$ and
$a_n/(2b_n)=1-\frac1n+O(n^{-2})$, in agreement with Theorem~\ref{thm:pure}.
\end{remark}

\begin{remark}[Hankel determinants of zeta values]\label{rem:pure-zeta}
% src: Paper 1 Corollary 1.4; novelty_G1RH.md §1.2.
For $a>0$ and $2a+b>1$ the measure $\sigma_{a,2a+b}$ has moments $\int x^j\diff\sigma_{a,2a+b}=\zeta(aj+2a+b)$, so
its Hankel determinants are $\det_{1\le i,j\le n}\zeta(a(i+j)+b)$ \cite[Corollary~1.4]{Paper1}. Theorem~\ref{thm:pure}
therefore gives $n^ab_n=c(a)(1+O(1/n))$ for the recurrence coefficients behind these determinants. For $a=1$ Monien
\cite{Monien} reports a full asymptotic expansion of the determinants, obtained numerically by D.~Zagier; it goes
beyond Theorem~\ref{thm:pure} but is not proved.
\end{remark}

\begin{remark}[where $c(p)$ comes from]\label{rem:pure-constant}
The equilibrium measure $\mu^*$ of \cite{Paper1} is supported on $(0,\infty)$ and coincides with the node density on
$[4c(p),\infty)$ (it is \emph{saturated} there); its unconstrained part is the band $(0,4c(p))$
(Proposition~\ref{prop:equilibrium}). After rescaling by $n^p$, the Riemann--Hilbert analysis sees a band of length
$4c(p)$, and a band of length $\beta$ produces off-diagonal entries close to $\beta/4$.
\end{remark}

The third result is elementary, and it is the reason why the analytic part of the proof only has to cover a window of
exponents $s$. A sequence $\eta_n>0$ is a \emph{regular rate} if $\eta_n$ is non-increasing, $n\eta_n$ is non-decreasing,
and $\eta_n\to0$; for example $\eta_n=1/n$.
% src: ROUTE2_alpha.md §4.

\begin{maintheorem}[Christoffel shift: law and rate]\label{thm:shift}
% src: PROOF.md §3.21 (Theorems SH and SH-r; Lemmas CP, CT, CS, CS-r); ROUTE2_alpha.md §§2-4.
Let $\sigma$ be a positive Borel measure on $[0,\infty)$ with infinite support and finite moments, let $p>0$ and
$c>0$, and let $\eta_n$ be a regular rate. If
\[
  |n^pa_n(\sigma)-2c|+|n^pb_n(\sigma)-c|=O(\eta_n),
\]
then the measure $x\sigma$ \textup{(}$\diff(x\sigma):=x\diff\sigma(x)$\textup{)} satisfies the same estimate, with
the same $c$ and rate $O(\eta_n)$. Without a rate: if $n^pa_n(\sigma)\to2c$ and $n^pb_n(\sigma)\to c$, then the same
holds for $x\sigma$. For $x_k=k^{-p}$ one has exactly $x\cdot\sigma_{p,s,\ell}=\sigma_{p,s+p,\ell}$.
\end{maintheorem}

\begin{remark}\label{rem:shift-credit}
The Jacobi matrix of $x\sigma$ is obtained from that of $\sigma$ by factorising $J=LU$ and reversing the factors; these
transfer formulas are classical \cite[\S1]{DD}. The factors are governed by a chain sequence, and the statement
without rate follows from a proposition of Dai, Piao and Qiao \cite[Proposition~2.1]{DPQ} together with those formulas.
The part with a rate is the one that matters here: it lets the rate $O(1/n)$ survive the shift. The parameter
sequence of the chain converges only at the rate $O(\sqrt{\eta_n})$ in general, and this is sharp; the coefficients of
$x\sigma$ involve only its increments, which are $O(\eta_n)$ (Section~\ref{sec:shift}). We have not found the rate
statement in the literature.
\end{remark}

\begin{remark}\label{rem:shift-converse}
% src: PROOF.md §3.19 Lemma CT_1 "<=" with Lemma D(b); novelty_lineB_A.md §2.
Conversely, without rate: if $x\sigma$ satisfies \eqref{eq:law} (with any $c>0$ in place of $c(p)$), so does
$\sigma$. In particular the pointwise law is insensitive to a point mass at $0$. This direction is also a consequence
of \cite[Proposition~2.1]{DPQ}, which applies to every parameter sequence of a chain sequence
(Section~\ref{sec:shift}).
\end{remark}

\subsection{Context}\label{sec:context}
% src: research/program/novelty_lineB_A.md, novelty_lineB_B.md, novelty_G1RH.md; RESEARCH_PROGRAM.md §§1, 6.
\emph{Averaged versus pointwise.} The averaged law \eqref{eq:paper1-averaged} is a statement about $\log D_n$, the
logarithm of the Hankel determinant, at order $n^2$, and it is proved in \cite{Paper1} by a discretised equilibrium
problem. Such arguments are insensitive to bounded perturbations of the weights: multiplying $w_k$ by factors in
$[e^{-\varepsilon},e^{\varepsilon}]$ changes $\log D_n$ by at most $\varepsilon n$. They therefore cannot yield
\eqref{eq:law}, which is sensitive to such perturbations (see the announcement below).

\emph{Riemann--Hilbert analysis of discrete orthogonal polynomials.} The Riemann--Hilbert (RH) approach to discrete
orthogonal polynomials was developed by Baik, Kriecherbauer, McLaughlin and Miller \cite{BKMM} for measures on finite
node sets in a bounded interval, with real-analytic node density and potential; their analysis includes saturated
regions, which also appear here. Within the literature we have read, the only Riemann--Hilbert analysis of discrete orthogonal
polynomials on infinitely many nodes accumulating at a finite point is that of Wu, Lin, Xu and Zhao \cite{WLXZ}, for
the Tricomi--Carlitz polynomials; our frame follows theirs. Dai and Yao make the same observation about the
literature \cite[\S1]{DY}.
% TODO(citation gate): Dai--Yao quoted only via novelty_G1RH.md §2; re-read before submission (OUTLINE §5, §1.3(b)).

\emph{Non-analytic weights.} Our weights are not analytic at the hard edge, and the standard local parametrix does not
exist there. The same difficulty was met by Conway and Deift \cite{CD} for the weight $\log\frac{2k}{1-x}$ on
$(-1,1)$, which carries a slowly varying factor at the endpoint $x=1$. They write that ``there is no known parametrix
for the Riemann--Hilbert problem in a neighborhood of the logarithmic singularity at $x=1$'', and they compare the
problem with the Legendre problem through the operator theory behind RH problems instead of through a quotient.
The comparison argument in Section~\ref{sec:hardedge} is that of Conway and Deift \cite[Theorem~4.7 and
Proposition~4.9]{CD}. What we add is a comparison model built from the Szeg\H{o} function of the full weight, a uniform
resolvent bound for a Bessel edge of any order $|\alpha|<1$, and the treatment of an arbitrary slowly varying factor.
Deift and Piorkowski \cite{DP} extended the analysis of \cite{CD} to the weight $\log\frac{2}{1-x}$ and note that, for a
comparison with a weight that has a simple zero at the endpoint, ``the crucial Theorem~4.7 in [4] requires a different
proof'' ([4] being \cite{CD}); this is consistent with our restriction to Bessel orders $|\alpha|<1$
(Remark~\ref{rem:alpha-ge-1}). Other approaches to non-analytic weights are Piorkowski's parametrix-free method, which
needs an a priori estimate for the exact solution \cite{Pio}, the $\bar\partial$ methods of McLaughlin and Miller
\cite{MM} and of Trogdon \cite{Tro}, which assume a finite number of Lipschitz derivatives of the weight (for
\cite{Tro}, $C^{k,1}$ regularity up to the endpoints of the support), and the
approximation of a smooth symbol by analytic ones in Deift, Its and Krasovsky \cite[\S6.2]{DIK}. Our reduction from
arbitrary to analytic slowly varying factors (Section~\ref{sec:smoothing}) is of the last kind, but the approximation
error is controlled outside the RH analysis.

\emph{Recurrence asymptotics without analyticity.} For exponential weights $W=e^{-Q}$, Levin and Lubinsky
\cite{LL} prove $o(1)$ asymptotics of the recurrence coefficients, with constants given by equilibrium quantities, under
regularity conditions that do not involve analyticity, and without RH methods. Theorem~\ref{thm:main} is an analogue
for compact Jacobi operators. We have not found Theorem~\ref{thm:main} in the literature, even for bounded $\ell$. The
closest results we know concern continuous weights on an interval \cite{CD,DP,Pio,Tro,MM} or exponential weights
\cite{LL}.

\emph{Special cases.} For $(p,s)=(2,2)$ the coefficients are explicit (Lommel polynomials,
Remark~\ref{rem:pure-lommel}). The measures $\sigma_{a,2a+b}$ produce the Hankel determinants of zeta values studied by
Monien \cite{Monien} and by Haynes and Zudilin \cite{HaynesZudilin} (Remark~\ref{rem:pure-zeta}). The orthogonality
measure of the Tricomi--Carlitz polynomials is supported on $\{\pm1/\sqrt{k+\alpha}\}$ \cite[\S3]{VA}; in the squared
variable its atoms are $1/(k+\alpha)$, so it is an instance of $p=1$ with weights that are not pure powers. For it the
recurrence coefficients are explicit \cite[\S3]{VA}, and \cite{WLXZ} obtain uniform asymptotics of the polynomials.

\emph{Rough weights.} Slow variation cannot be replaced by boundedness. In a further paper \cite{Paper3} we show the
following. For every $p>0$ and $s>1$ there is $\varepsilon_0>0$ such that, for every $\varepsilon\in(0,\varepsilon_0]$,
there are weights $w_k=k^{-s}e^{f(k)}$ with $|f(k)|\le\varepsilon$ for which $n^pb_n$ does not converge, and weights of
the same kind for which $n^pa_n$ does not converge. The averaged law \eqref{eq:paper1-averaged} holds for all of them.
Together with Theorem~\ref{thm:main} this gives a dichotomy: the pointwise law is stable under slowly varying
perturbations of the weights, of any size, and unstable under rough ones, however small.
% src: THEOREMS.md H4 Theorem R (b-law), §2.6 Theorem Ra (a-law). No constants; not "the same weights break both".

\subsection{Strategy of the proof}\label{sec:strategy}
The proof has a soft part and a hard part. The soft part (Sections~\ref{sec:shift} and~\ref{sec:smoothing}) is
elementary and self-contained; it reduces Theorem~\ref{thm:main} to the following statement, which is the technical
core of the paper.

\begin{maintheorem}[the Riemann--Hilbert engine: analytic factor, base window]\label{thm:engine}
% src: ROUTE2.md §4.4 Theorem RII (p > 1); ROUTE2_xframe1.md §4 Theorem RII_x1 (p = 1, 1 < r < 3, incl. alpha = 0);
% ROUTE2_plt1.md §4.3 Theorem RII_{<1} (0 < p < 1).
Let $p>0$ and $1<s<1+2p$, i.e.\ $\alpha:=(s-1-p)/p\in(-1,1)$. Let $\ell_1$ satisfy \Ath{\vartheta_p}. Then
$\sigma_{p,s,\ell_1}$ satisfies
\[
  n^pb_n=c(p)\bigl(1+O(1/n)\bigr),\qquad a_n/(2b_n)=1+O(1/n).
\]
\end{maintheorem}

\begin{remark}[what is, and is not, proved about the slowly varying part]\label{rem:engine-scope}
% src: ROUTE2 Theorem RII, "What is proved ..." (REVIEW_Plan48 F3); OUTLINE §7 R2.
The proof compares the true problem with an exactly solvable comparison problem. Relative to the comparison problem,
the non-analytic part of $\ell_1$ moves the first two coefficients at infinity by $O(\varepsilon'_n/n)$, where
$\varepsilon'_n\to0$ (Section~\ref{sec:hardedge}). The comparison problem itself depends on $\ell_1$ at order $1/n$.
So no claim is made that $b_n(\sigma_{p,s,\ell_1})/b_n(\sigma_{p,s})-1=O(\varepsilon_\Lambda/n)$; the $O(1/n)$ of
Theorem~\ref{thm:engine} is absolute, not relative.
\end{remark}

\emph{From Theorem~\ref{thm:engine} to Theorems~\ref{thm:main} and~\ref{thm:pure}.} Write $s=s_0+mp$ with
$s_0\in(1,1+2p)$ and an integer $m\ge0$; this is possible because the window $(1,1+2p)$ has length $2p>p$.
Theorem~\ref{thm:engine} at $s_0$ and $m$ applications of Theorem~\ref{thm:shift} with $\eta_n=1/n$ give the law with
rate $O(1/n)$ for every analytic factor and every $s>1$. This is Theorem~\ref{thm:main}(iii), and with $\ell\equiv1$ it
is Theorem~\ref{thm:pure}. For a general slowly varying $\ell$ we apply this to the analytic factor
$\ell_1:=\exp\Lambda_\ell(\log\cdot)$ and compare: a multiplicative perturbation $\ell/\ell_1\to1$ of the weights
changes $\log b_n^2$ and $n^pa_n$ by at most a constant times the supremum of $|\log(\ell/\ell_1)|$ over the indices
$k\ge n/M_0$, where the spectral weight of $P_n$ lives, up to an exponentially small contribution from the head
$k<n/M_0$ (Proposition~\ref{prop:stability}). This gives Theorem~\ref{thm:main}(ii), and (i) follows.

\emph{Why the window $|\alpha|<1$.} Near the accumulation point $0$ the measure looks like a continuous weight: the
nodes $x_k=k^{-p}$ have density $\frac1px^{-1-1/p}$, and the mass per unit length near $x$ is
$k^{-s}\cdot\frac1px^{-1-1/p}=\frac1px^{\alpha}$ with $x=k^{-p}$. In the RH analysis $0$ becomes a hard edge with Bessel
order $\alpha=(s-1-p)/p$. Our comparison argument works for $|\alpha|<1$ and fails for $\alpha\ge1$
(Remark~\ref{rem:alpha-ge-1}). Multiplication by $x$ raises $\alpha$ by $1$, and Theorem~\ref{thm:shift} carries the
law and its rate across. For $s>1+p/2$ one can even choose $s_0$ with $|\alpha|\le\frac12$; the case $\alpha=0$ cannot be
avoided in general (Remark~\ref{rem:base-window}).

\emph{The frame.} After rescaling $z\mapsto n^pz$ the nodes $(n/k)^p$ fill $(0,n^p]$. The equilibrium measure $\mu^*$
of \cite{Paper1} governs the $n$-th root behaviour: on the band $(0,4c(p))$ the polynomials oscillate, on the saturated
region $(4c(p),n^p)$ the zero distribution coincides with the node density, and $4c(p)$ is a soft edge with an
Airy parametrix. Near $0$ the band meets infinitely many nodes. We set up the discrete interpolation problem and its
first transformations as Wu, Lin, Xu and Zhao do \cite{WLXZ}, in the one-sided variable $x$ (Section~\ref{sec:frame}).
What depends on $p$ is the outer scalar function that removes the Gamma-function growth of the node masses: the index
function $t(z)=nz^{-1/p}$ maps the slit plane onto the sector $|\arg t|<\pi/p$, which lies inside, exactly fills, or
exceeds the sector $|\arg t|<\pi$ where Stirling's formula holds, according as $p>1$, $p=1$ or $p<1$.
Section~\ref{sec:outer} treats the three cases in turn and ends with a Frame Proposition (Proposition~\ref{prop:frame}),
which is the only interface between the frame and the hard edge.

\emph{The hard edge, and why a local parametrix fails.} The slowly varying factor enters the band weight through
$r_n(z):=\ell_1(nz^{-1/p})/\ell_1(n)$. This function is analytic off $(-\infty,0]$, but it is not analytic at $0$, and it
is not of Fisher--Hartwig type. The standard Bessel parametrix needs a square root of the band weight continued across
$(-\delta,0)$. Any such continuation jumps across $(-\delta,0)$ by a factor $1+O(\varepsilon_\Lambda)$, with no rate in
$n$. Conjugated by the Bessel model at the scale $|z|\asymp n^{-2}$, this jump becomes of size
$O(n\varepsilon_\Lambda)$, which is not small. A Bessel order that drifts with $n$ removes only the leading,
power-like part of the factor. So a quotient (small-norm) argument fails.
% src: ROUTE2 §2, last paragraph.

\emph{The remedy.} Following Conway and Deift, we build the Bessel parametrix from the Szeg\H{o} function $D_W$ of the
whole band weight, the slowly varying factor included. This function has no jump on $(-\delta,0)$, so the comparison
problem is exactly solvable. The error moves to the lips of the lens, where the jump of the true problem differs from
that of the comparison problem by an amount that is small in the supremum norm but has no rate. Such smallness is
exactly what a Neumann series around the comparison resolvent needs, once that resolvent is bounded uniformly in $n$
(Proposition~\ref{prop:UR-comparison}; this is where $|\alpha|<1$ is used, Appendix~\ref{app:resolvent}). Finally
$e^{2n\phi}$ confines the mismatch to $|z|\lesssim n^{-2}$, up to a tail, and this Bessel localisation supplies the
factor $1/n$.

\emph{What is new relative to Conway--Deift.} The weight factor is an arbitrary analytic slowly varying function rather
than one explicit logarithm; the Bessel order is any $\alpha\in(-1,1)$ rather than the Legendre value $\alpha=0$; the
comparison model is an explicit parametrix rather than a second orthogonality problem; and the problem is discrete,
with nodes accumulating at the hard edge. The key estimate for the slowly varying factor is a scale-invariant path
integral of $\Lambda'$ (Section~\ref{sec:hardedge}).

Figure~\ref{fig:dependency} shows how the parts fit together.

\begin{figure}[ht]
\centering
\begin{tikzpicture}[
  box/.style={draw,rounded corners=2pt,align=center,font=\footnotesize,inner sep=3pt,text width=3.5cm,
    minimum height=1.05cm},
  res/.style={box,fill=black!6,thick},
  arr/.style={-{Stealth[length=2mm]}},
  lab/.style={font=\scriptsize,align=center}]
  \node[box] (P1) at (0,0) {Paper 1: equilibrium measure $\mu^*$, band $[0,4c(p)]$};
  \node[box] (FR) at (0,-1.6) {\S\S5--6: frame, outer scalar in three geometries, Frame Proposition};
  \node[box] (HE) at (0,-3.2) {\S7: hard edge; comparison with a Bessel model, $|\alpha|<1$};
  \node[res] (T4) at (0,-4.8) {\S8: Theorem~\ref{thm:engine}\\ (analytic $\ell_1$, $1<s<1+2p$)};
  \node[box] (TL) at (5.9,-1.6) {Tools: discrete RH frame \cite{WLXZ}, Bessel model \cite{KMVV}, comparison \cite{CD}};
  \node[box] (AN) at (5.9,-4.8) {analytic $\ell_1$, every $s>1$,\\ rate $O(1/n)$};
  \node[res] (T2) at (11.8,-3.2) {Theorem~\ref{thm:pure}\\ ($\ell\equiv1$)};
  \node[res] (T1) at (11.8,-4.8) {Theorem~\ref{thm:main}\\ (any slowly varying $\ell$)};
  \draw[arr] (P1) -- (FR);
  \draw[arr] (FR) -- (HE);
  \draw[arr] (HE) -- (T4);
  \draw[arr] (TL) -- (FR);
  \draw[arr] (TL.south) |- (HE.east);
  \draw[arr] (T4) -- node[lab,above]{\S3: shift} node[lab,below]{Thm.~\ref{thm:shift}} (AN);
  \draw[arr] (AN) -- node[lab,above]{\S4: smoothing} node[lab,below]{and stability} (T1);
  \draw[arr] (AN.north east) -- node[lab,above left]{$\ell\equiv1$} (T2.west);
\end{tikzpicture}
\caption{Structure of the proof. Sections~\ref{sec:shift} and~\ref{sec:smoothing} are elementary and reduce
Theorems~\ref{thm:main} and~\ref{thm:pure} to Theorem~\ref{thm:engine}, which is proved in
Sections~\ref{sec:frame}--\ref{sec:extraction}.}
\label{fig:dependency}
\end{figure}

\subsection{Organisation}\label{sec:organisation}
Section~\ref{sec:prelim} fixes conventions, proves two facts about slowly varying sequences, quotes the input from
\cite{Paper1} (only parts (i) and (ii) of its Theorem~1.3 are used), and states the external tools with their
hypotheses. Section~\ref{sec:shift} proves Theorem~\ref{thm:shift}. Section~\ref{sec:smoothing} reduces arbitrary
slowly varying factors to analytic ones and proves Theorems~\ref{thm:main} and~\ref{thm:pure}, assuming
Theorem~\ref{thm:engine}. Sections~\ref{sec:frame}--\ref{sec:extraction} prove Theorem~\ref{thm:engine}: the common
frame (Section~\ref{sec:frame}), the outer scalar function in the three geometries $p>1$, $p=1$ and $p<1$
(Section~\ref{sec:outer}), the comparison at the hard edge (Section~\ref{sec:hardedge}), and the extraction of $a_n$ and
$b_n$ (Section~\ref{sec:extraction}). Section~\ref{sec:numerics} illustrates the results numerically, and
Section~\ref{sec:further} collects remarks and open problems. The appendices contain two $L^2$ bounds on a star of arcs
(Appendix~\ref{app:L2}), the proof of the uniform resolvent bound (Appendix~\ref{app:resolvent}) and the a priori bound
at the hard edge (Appendix~\ref{app:apriori}). A reader who is mainly interested in the Riemann--Hilbert analysis can
start at Section~\ref{sec:frame}, with the statement of Theorem~\ref{thm:engine} in hand.

Throughout, $p>0$ is fixed, $C,C'$ denote positive constants that may change from line to line (never $c$, which is
reserved for $c(p)$), and $n$ is the degree.
